%% kapa: exam=1 round=30 year=2019 period=2 form=A subject=accounting subject_ko=회계학 seq=18 number=58
%% source: Q-Net 감정평가사 기출문제 (cst00309) artlSeq=5205205 file=extracted/1008/5205205/2019년 제30회 감정평가사 1차시험 문제지/2019년 제30회 감정평가사 1차시험 문제지/2019년도 제30회 감정평가사 1차 2교시 A형.hwp
%% conversion: kordoc 4.24.51 -> tools/gen1.py (auto); review=proofread-sonnet; flags=none
\begin{question}{58}
재무제표 요소의 측정에 관한 설명으로 옳지 않은 것은?
\begin{choices}
  \item 재무제표를 작성할 때 기업이 가장 보편적으로 채택하고 있는 재무제표 요소의 측정기준은 역사적 원가이다.
  \item 재무제표를 작성할 때 합리적 추정을 사용해야 하는데 이는 신뢰성을 훼손하게 된다.
  \item 부채에 현행원가 개념을 적용하면 현재시점에서 그 의무를 이행하는데 필요한 현금이나 현금성자산의 할인하지 아니한 금액으로 평가한다.
  \item 자산에 대하여 손상차손회계를 적용할 때 고려하는 사용가치는 그 자산의 공정가치와 다르다.
  \item 재무제표 요소의 측정은 재무상태표와 포괄손익계산서에 인식되고 평가되어야 할 재무제표 요소의 화폐금액을 결정하는 과정으로 특정 측정기준의 선택과정을 포함한다.
\end{choices}
\end{question}
\answer{2}
